%%=============================================================================
%% RESULTADOS
%%=============================================================================
\chapter{Resultados}
\label{cap:resultados}

Este capítulo presenta los resultados de los seis experimentos descritos en la Metodología (Capítulo~\ref{cap:metodologia}). Cada sección corresponde a un experimento. Salvo que se indique lo contrario, las estimaciones de min-entropía provienen de las herramientas oficiales del NIST (\texttt{ea\_non\_iid} y \texttt{ea\_iid}) y las pruebas de SP 800-22, de STS 2.1.2.

\section{E1: la referencia es determinística en microcontroladores}
\label{sec:res_determinismo}

La Tabla~\ref{tab:res_determinismo} resume las capturas de la traza $(n_p, \Delta t)$ de la referencia tras varios reinicios de cada placa, con la misma semilla.

\begin{table}[H]
\centering
\caption{Repetibilidad de la referencia entre reinicios}
\label{tab:res_determinismo}
\begin{tabularx}{\textwidth}{@{}Ycccrr@{}}
\toprule
\textbf{Plataforma} & \textbf{Reinicios} & \textbf{Registros} & \textbf{¿Idénticas?} & \textbf{Ciclos/\textit{pad}} & \textbf{B/s} \\
\midrule
Mega 2560 & 3 & 200--300 & sí, 100\,\% & 4\,374 & 5.95 \\
Nano & 2 & 300 & sí, 100\,\% & 4\,212 & 6.15 \\
Nano, referencia con \texttt{opt-level=\textquotedbl{}s\textquotedbl{}} & 2 & 300 & sí, 100\,\% & 4\,259 & 5.87 \\
nRF52840 & 3 & 2\,000 & sí, 100\,\% & 364 & 287 \\
Apple Silicon (control) & 3 & 64 B & no & -- & -- \\
\bottomrule
\end{tabularx}
\end{table}

En los tres microcontroladores, todas las capturas fueron idénticas bit a bit. La razón es que $\Delta t$ es una función determinística de los valores del generador interno (el número de \textit{pads}, los rechazos de la reducción de Lemire y las interrupciones de desbordamiento del temporizador), todos ellos funciones de la semilla, y la siguiente semilla es a su vez una función de $\Delta t$: no hay ruido independiente en ninguna parte del ciclo. La traza del Nano difiere de la del Mega (otra disposición del código produce otros conteos de ciclos), pero cada placa se repite a sí misma con exactitud. Compilar la referencia con optimización por tamaño en lugar de \texttt{-O0} no cambia el determinismo ni, de forma notable, la velocidad: en un núcleo de 8 bits domina la aritmética de 64 y 128 bits.

Las pruebas estadísticas no detectan esta falla. El byte bajo de $\Delta t$ en la traza del nRF52840 obtiene 4.93 bits por muestra con los estimadores MCV y \textit{t-tuple}, y la salida por defecto reconstruida a partir de la traza supera todas las pruebas de nivel 1 aplicables a 400 bytes. Solamente la comparación entre reinicios revela que la fuente no tiene entropía: la referencia, en estas plataformas, se comporta como un generador pseudoaleatorio cuya clave es la semilla fijada en tiempo de compilación.

\section{E2: fuentes físicas independientes}
\label{sec:res_fuentes}

La Tabla~\ref{tab:res_fuentes} resume la caracterización de las fuentes físicas. Las capturas de 20\,000 muestras tomaron unos 5.5 minutos por AVR y 35 segundos en el nRF52840; las de 1\,000\,001 muestras, del orden de 4.7 horas por AVR y 29 minutos en el nRF52840.

\begin{table}[H]
\centering
\caption{Caracterización de las fuentes físicas (min-entropía por muestra según \texttt{ea\_non\_iid})}
\label{tab:res_fuentes}
\small
\begin{tabularx}{\textwidth}{@{}Y R{0.16\textwidth} R{0.05\textwidth} R{0.10\textwidth} R{0.08\textwidth} R{0.07\textwidth} R{0.12\textwidth}@{}}
\toprule
\textbf{Fuente} & \textbf{Periodo medio (ciclos)} & \textbf{$\sigma$} & \textbf{Autocorr. lag 1} & \textbf{$H$ ($2 \cdot 10^4$)} & \textbf{$H$ ($10^6$)} & \textbf{Acreditado} \\
\midrule
Mega 2560, perro guardián & 271\,514 (16.97~ms) & 18.3 & 0.65 & 2.14 & \textbf{1.62} & 1 \\
Nano, perro guardián & 265\,950 (16.62~ms) & 45.1 & 0.69 & 2.98 & \textbf{2.24} & 1 \\
nRF52840, LFRC $\times 32$ & 111\,514 & 37.0 & $-0.91$ & 0.61 & \textbf{0.60} & no se usa \\
nRF52840, RNG en hardware & -- & -- & -- & 6.89 & \textbf{7.24} & 4 \\
\bottomrule
\end{tabularx}
\end{table}

En ambos AVR, el estimador determinante fue \textit{t-tuple}. Con un millón de muestras, ese estimador encuentra más estructura que con 20\,000, porque el periodo del perro guardián deriva lentamente (autocorrelación de 0.65 a 0.69), y los valores estimados bajaron entre 24 y 25\,\%. La Tabla~\ref{tab:res_estimadores_avr} detalla los estimadores. El valor acreditado de 1 bit por muestra sigue cubierto, con un margen de $1.6\times$ en el Mega y de $2.2\times$ en el Nano. Con 8 muestras por byte, cada byte de salida queda respaldado por al menos 12.9 bits estimados en el Mega y 17.9 en el Nano.

\begin{table}[H]
\centering
\caption{Estimadores de \texttt{ea\_non\_iid} sobre $10^6$ muestras del perro guardián (bits por muestra)}
\label{tab:res_estimadores_avr}
\begin{tabular}{@{}lcccccccc@{}}
\toprule
\textbf{Fuente} & \textbf{MCV} & \textbf{\textit{t-tuple}} & \textbf{LRS} & \textbf{MultiMCW} & \textbf{Lag} & \textbf{MultiMMC} & \textbf{LZ78Y} & \textbf{Estimado} \\
\midrule
Mega 2560 & 4.56 & \textbf{1.62} & 1.86 & 3.30 & 2.01 & 2.31 & 3.78 & \textbf{1.62} \\
Nano & 4.91 & \textbf{2.24} & 3.38 & 2.93 & 3.30 & 2.92 & 2.93 & \textbf{2.24} \\
\bottomrule
\end{tabular}
\end{table}

La fuente LFRC del nRF52840 presentó solo 42 valores distintos y una fuerte correlación negativa entre muestras consecutivas. El oscilador de 32.768 kHz es mucho más estable que el del perro guardián de AVR, y lo que varía es principalmente la cuantización de la lectura del contador del RTC: una marca de tiempo que llega tarde en un flanco llega temprano en el siguiente. Como el chip tiene un generador en hardware, esta fuente se descartó.

El periférico RNG del nRF52840, con un millón de muestras, obtuvo 7.24 bits por byte con \texttt{ea\_non\_iid} (MCV 7.88, \textit{t-tuple} 7.35, LRS 7.72, predictores entre 7.91 y 7.95; el valor determinante proviene de la evaluación por bits, multiplicada por 8) y superó todas las pruebas IID de \texttt{ea\_iid} (independencia por $\chi^2$, $p = 0.94$). Se le acreditaron 4 bits por byte, el mayor valor que cubre la tabla de umbrales de las pruebas de salud, con un margen de más de 3 bits.

Una segunda captura independiente de cada fuente difirió de la primera desde la primera muestra; para el RNG del nRF52840, hubo 79 coincidencias posicionales en 20\,000 muestras, frente a 78 esperadas por azar. Esta es exactamente la propiedad de la que carece la referencia. Como era de esperar, las muestras crudas fallan gravemente las pruebas de nivel 1 ($\chi^2 \approx 10^5$, correlación serial entre 0.4 y 0.65): son una fuente de ruido con deriva lenta, no bytes uniformes, y por eso el generador las procesa y acondiciona.

\section{E3: la variante IoT sobre hardware}
\label{sec:res_iot}

\subsection{Optimización del cómputo}

La Tabla~\ref{tab:res_optimizacion} muestra la evolución del costo de cómputo por byte en el ATmega2560, medida con una fuente sintética para aislar el cómputo de la espera por muestras.

\begin{table}[H]
\centering
\caption{Costo de cómputo por byte de salida en el ATmega2560}
\label{tab:res_optimizacion}
\begin{tabularx}{\textwidth}{@{}Xrr@{}}
\toprule
\textbf{Paso} & \textbf{Ciclos por byte} & \textbf{Mejora} \\
\midrule
Primera versión ($N = 6$, 8 caminatas por byte, $256 \bmod b$ por índice) & 10\,733\,218 & -- \\
Umbral de Lemire tomado de una tabla (AVR no tiene divisor) & 4\,585\,746 & $2.3\times$ \\
Generación y composición fusionadas, salida temprana & 2\,863\,887 & $1.6\times$ \\
4 caminatas por byte, 2 muestras por caminata & \textbf{1\,621\,365} & $1.8\times$ \\
\bottomrule
\end{tabularx}
\end{table}

En total, el cómputo por byte se redujo $6.6\times$. Por índice de Fisher--Yates, el costo pasó de unos 1\,090 ciclos en la referencia a unos 113 en la variante IoT ($9.7\times$), y por \textit{pad} se redujo cerca de $7.8\times$, aunque cada \textit{pad} de la variante IoT tiene un índice más ($N = 6$). Con 1.62 millones de ciclos (101 ms) por byte frente a 8 muestras del perro guardián por byte ($8 \times 17 = 136$ ms), el generador en AVR quedó limitado por la fuente y no por el cómputo: la caminata termina mientras la interrupción del perro guardián todavía está acumulando las siguientes muestras.

\subsection{Divergencia, rendimiento y huella}

La Tabla~\ref{tab:res_iot_hw} resume la ejecución de la variante IoT en las tres placas. En los tres casos, la salida difirió entre reinicios y no hubo fallas de las pruebas de salud.

\begin{table}[H]
\centering
\caption{Variante IoT sobre hardware}
\label{tab:res_iot_hw}
\begin{tabularx}{\textwidth}{@{}L{0.15\textwidth}>{\raggedright\arraybackslash}XL{0.17\textwidth}rcr@{}}
\toprule
\textbf{Placa} & \textbf{Entre reinicios (coincidencias)} & \textbf{Bytes (tiempo)} & \textbf{B/s} & \textbf{Fallas de salud} & \textbf{Flash} \\
\midrule
Mega 2560 & distintas, 0/40 & 10\,000 (22.6 min) & 7.36 & 0 & 6\,470 B \\
Nano & distintas, 0/40 & 10\,000 (22.2 min) & 7.52 & 0 & 6\,330 B \\
nRF52840 & distintas, 0/200 & 50\,000 (8.2 min) & 101 & 0 & 2\,940 B \\
\bottomrule
\end{tabularx}
\end{table}

Los tamaños corresponden al \textit{firmware} completo, que incluye el código de la UART, la medición de cómputo al arranque y dos instancias del generador (con la fuente real y con la sintética). El \textit{firmware} de la referencia, compilado con \texttt{-O0}, ocupa 8\,912, 8\,780 y 3\,292 bytes, respectivamente.

La Tabla~\ref{tab:res_rendimiento} compara el rendimiento de ambas versiones. En AVR, la variante IoT es más rápida que la referencia y, a diferencia de esta, cada byte está respaldado por al menos 8 bits acreditados. En el nRF52840 la variante IoT es más lenta, porque las caminatas con $N = 6$ son seis veces más largas que con $N = 5$, aunque por \textit{pad} sigue siendo $1.5\times$ más barata (236 frente a 364 ciclos).

\begin{table}[H]
\centering
\caption{Rendimiento de la salida (bytes por segundo)}
\label{tab:res_rendimiento}
\begin{tabularx}{\textwidth}{@{}L{0.15\textwidth}L{0.22\textwidth}L{0.22\textwidth}Y@{}}
\toprule
\textbf{Placa} & \textbf{Referencia (0 bits de entropía)} & \textbf{Variante IoT (con respaldo de entropía)} & \textbf{Límite de la variante IoT} \\
\midrule
Mega 2560 & 5.95 & 7.36 & Tasa del perro guardián \\
Nano & 6.15 & 7.52 & Tasa del perro guardián \\
nRF52840 & 287 & 101 & Cómputo: 680\,000 ciclos (10.6 ms) por byte \\
\bottomrule
\end{tabularx}
\end{table}

\subsection{Pruebas estadísticas sobre muestras cortas}

La Tabla~\ref{tab:res_iot_tier1} presenta las pruebas de nivel 1 y ENT sobre la salida. Todas las pruebas se superan en las tres placas. Como comparación, los bytes crudos del RNG del nRF52840 fallan la prueba de rachas ($p = 0.006$), mientras que, después de la caminata de QPP, la salida la supera.

\begin{table}[H]
\centering
\caption{Pruebas de nivel 1 y ENT sobre la salida de la variante IoT}
\label{tab:res_iot_tier1}
\small
\begin{tabularx}{\textwidth}{@{}Ycccccc@{}}
\toprule
\textbf{Flujo} & \textbf{Monobit $p$} & \textbf{Rachas $p$} & \textbf{$\chi^2$ ($p$)} & \textbf{Corr.\ serial} & \textbf{Shannon} & \textbf{Error de $\pi$} \\
\midrule
Mega 2560 (10 kB) & 0.91 & 0.45 & 277.4 (0.16) & $-0.0025$ & 7.980 & 0.27\,\% \\
Nano (10 kB) & 0.87 & 0.99 & 243.5 (0.69) & $-0.0020$ & 7.982 & 0.27\,\% \\
nRF52840 (50 kB) & 0.17 & 0.80 & 239.8 (0.74) & $-0.0010$ & 7.997 & 0.45\,\% \\
\textit{RNG crudo del nRF52840 (20 kB)} & 0.26 & \textbf{0.006} & 294.6 (0.045) & $-0.0005$ & 7.989 & 0.74\,\% \\
\bottomrule
\end{tabularx}
\end{table}

En la vía IID de SP 800-90B, las 19 pruebas de permutación y la prueba LRS se superaron en todas las salidas. La prueba de independencia por $\chi^2$ falló en los dos AVR (estadístico de 951 y 946 con 735 grados de libertad), pero falló de la misma manera en tres muestras de 10 kB de \texttt{/dev/urandom} y en dos fragmentos de 10 kB de la salida del nRF52840, que sí la supera con 50 kB. Con 256 símbolos y 10\,000 muestras, el conteo esperado por par de símbolos es de aproximadamente 0.15, de modo que la prueba agrupa casi todos los pares en pocas clases y su valor $p$ no está calibrado (SP 800-90B especifica al menos un millón de muestras). Una verificación directa sobre las salidas AVR (pares de nibbles altos y bajos y de los bits 0 y 7, con desfases de 1 a 5 y 8) no encontró dependencia consistente. Por lo tanto, se trata de un artefacto del tamaño de muestra y no de una propiedad del generador.

En la vía no IID, \texttt{ea\_non\_iid} estimó 5.81 bits por byte para el Mega, 5.75 para el Nano y 6.47 para el nRF52840. Estas cifras deben leerse junto a los controles: doce muestras de 10 kB de \texttt{/dev/urandom} obtuvieron entre 5.69 y 6.64, y una de 50 kB, 6.48. Con estos tamaños de muestra, el estimador no puede asignar más de unos 6.5 bits por byte ni siquiera a una fuente perfecta, y las salidas AVR quedan dentro de la dispersión de \texttt{/dev/urandom}. Las cifras no indican nada en contra del generador, pero tampoco pueden confirmar 8 bits por byte: para la salida, la afirmación de entropía se apoya en la evaluación de la fuente cruda de la Sección~\ref{sec:res_fuentes}, como prevé SP 800-90B.

\subsection{Batería completa sobre 1\,000\,000 de bytes}

Se capturó 1\,000\,000 de bytes de salida de la variante IoT en cada placa, con el mismo \textit{firmware} y los mismos valores acreditados de las secciones anteriores. En el nRF52840, la captura tomó 2.75 horas a 101 B/s. En los AVR, se dividió en 13 bloques (12 de 80\,000 bytes y uno de 40\,000), cada uno iniciado con un reinicio de la placa, y tomó 37.7 horas en el Mega 2560 (7.37 B/s) y 36.9 horas en el Nano (7.53 B/s). Ningún bloque presentó errores de sincronización ni fallas de las pruebas de salud. Los 13 bloques de cada placa son independientes entre sí: en los primeros 4\,000 bytes de cada uno de los 78 pares de bloques, la tasa de coincidencias posicionales fue de 0.398\,\% en el Mega y de 0.394\,\% en el Nano, frente al $1/256$ (0.391\,\%) esperado por azar, y ningún par compartió siquiera los primeros 4 bytes. La Tabla~\ref{tab:res_iot_1mb} presenta la batería completa sobre las tres salidas.

\begin{table}[H]
\centering
\caption{Batería completa sobre 1\,000\,000 de bytes de la variante IoT}
\label{tab:res_iot_1mb}
\small
\begin{tabularx}{\textwidth}{@{}L{0.20\textwidth}YYY@{}}
\toprule
\textbf{Prueba} & \textbf{Mega 2560} & \textbf{Nano} & \textbf{nRF52840} \\
\midrule
Nivel 1 & Todas superadas: monobit $p = 0.80$, rachas $p = 0.14$, $\chi^2$ $p = 0.67$, corr.\ serial $-0.0012$ & Todas superadas: monobit $p = 0.015$, rachas $p = 0.23$, $\chi^2$ $p = 0.23$, corr.\ serial 0.0017 & Todas superadas: monobit $p = 0.37$, rachas $p = 0.61$, $\chi^2$ $p = 0.50$, corr.\ serial 0.0009 \\
ENT (bits por byte; $\chi^2$ superado; error de $\pi$) & 7.999823; 66.58\,\%; 0.11\,\% & 7.999804; 23.31\,\%; 0.01\,\% & 7.999817; 49.98\,\%; 0.02\,\% \\
\texttt{ea\_iid} & Todas las pruebas IID superadas; $H = 7.889$ & Todas las pruebas IID superadas; $H = 7.870$ & Todas las pruebas IID superadas; $H = 7.866$ \\
\texttt{ea\_non\_iid} (bits por byte) & 7.358 & 7.261 & 7.277 \\
SP 800-22 (8 flujos de 1 Mbit) & 1\,419 de 1\,426 (99.51\,\%) & 1\,392 de 1\,400 (99.43\,\%) & 1\,440 de 1\,452 (99.17\,\%) \\
\bottomrule
\end{tabularx}
\end{table}

Las tres salidas superan todas las pruebas. En SP 800-22, el total de combinaciones varía porque las pruebas de excursiones aleatorias solo se aplican a los flujos con ciclos suficientes (5 de 8 en el Mega, 4 en el Nano y 6 en el nRF52840). Con $\alpha = 0.01$, se esperan por azar unas 14 fallas por placa, y se observaron 7, 8 y 12. Cada falla corresponde a un único flujo de los 8 que falla una subprueba: en el Mega, 6 plantillas de \textit{NonOverlappingTemplate} y la prueba \textit{Serial}; en el Nano, 7 plantillas y la prueba \textit{Universal}; en el nRF52840, 10 plantillas y 2 variantes de \textit{RandomExcursions}. Ninguna subprueba falla en más de un flujo. El valor de monobit del Nano ($p = 0.015$) supera el umbral de 0.01, y los 8 flujos del Nano superan la prueba \textit{Frequency} de SP 800-22, que mide la misma propiedad.

En \texttt{ea\_non\_iid}, el valor determinante en las tres placas proviene de la evaluación por bits, multiplicada por 8 (en el Mega, la estimación por símbolos fue de 7.889; en el Nano, de 7.334 por \textit{t-tuple}). Las tres cifras quedan entre 7.26 y 7.36 bits por byte, del mismo orden que la del RNG crudo del nRF52840 (7.24). En el nRF52840, el valor de la salida (7.28) es cercano al del RNG crudo, lo que indica que la caminata de QPP no añadió ni destruyó sesgo medible. En los AVR, la fuente cruda tiene solo entre 1.6 y 2.2 bits por muestra; las 8 muestras por byte respaldan al menos 12.9 bits evaluados por byte en el Mega y 17.9 en el Nano, de los cuales se acreditan 8, y la salida no muestra el sesgo de la fuente. Como en las muestras cortas, estas cifras son compatibles con una salida de entropía completa, pero no la demuestran: la afirmación de entropía de la variante IoT descansa en la evaluación de la fuente cruda de la Sección~\ref{sec:res_fuentes} y en las pruebas de salud continuas, no en la salida.

\section{E4: la referencia en anfitriones rápidos}
\label{sec:res_anfitriones}

La Tabla~\ref{tab:res_anfitriones} describe los dos anfitriones. En ambos, la calibración de la resolución del temporizador produjo $k = 1$. En la Raspberry Pi, la suite completa del espacio de trabajo pasó (118 pruebas superadas, ninguna fallida, una ignorada).

\begin{table}[H]
\centering
\caption{Anfitriones utilizados en E4}
\label{tab:res_anfitriones}
\begin{tabularx}{\textwidth}{@{}L{0.25\textwidth}XY@{}}
\toprule
& \textbf{Raspberry Pi 4 Model B} & \textbf{Computador de escritorio} \\
\midrule
CPU y sistema & Cortex-A72 ×4 (0.6--1.5 GHz), Debian 13 & Intel Core Ultra 7 265KF, Windows 11 \\
Temporizador & \texttt{CNTVCT\_EL0}, 54 MHz (18.5 ns por tick) & \texttt{QueryPerformanceCounter}, 10 MHz \\
$k$ nominal $\rightarrow$ calibrado & 34 $\rightarrow$ 1 & 100 $\rightarrow$ 1 \\
Ciclos por segundo & 349\,000 & 921\,000 \\
Salida por defecto & 69.8 kB/s & 184 kB/s \\
\bottomrule
\end{tabularx}
\end{table}

\paragraph{La referencia no es determinística en estos anfitriones.} Tres capturas de 20\,000 ciclos en la Pi empiezan con el mismo $n_p = 211$, determinado por la semilla, pero con $\Delta t$ de 303, 292 y 293 ticks, y divergen desde ahí. Las coincidencias posicionales de $n_p$ entre capturas fueron entre 73 y 89 de 20\,000, frente a unas 83 esperadas por azar. El modo de falla de E1 no se presenta.

\paragraph{El sesgo estructural se reproduce exactamente.} La salida de 1 MB reconstruida a partir de la traza falla las pruebas monobit, de rachas y de $\chi^2$ ($p \leq 10^{-24}$). El bit 7 vale 1 con frecuencia 0.48621, frente a 0.48654 calculado con la Ecuación~\eqref{eq:sesgo_xor} ($\sigma \approx 0.0005$); los bits 0 a 5 están en 0.5 dentro del ruido. El fallo es, entonces, el sesgo de la Sección~\ref{sec:sesgo}, y no un efecto de la plataforma.

\paragraph{Entropía fresca por ciclo.} En cada ciclo, la semilla absorbe $\Delta t \bmod 256$, pero $n_p$ ya está fijado por la semilla, de modo que la entrada nueva es $\Delta t$ \emph{dado} $n_p$. En la Pi, $\Delta t \approx 1.10 + 1.262\, n_p$ ticks: el 98\,\% de los ciclos cae entre $-1.5$ y $+0.7$ ticks de esa recta, y el percentil 99.9 está en $+4.7$ ticks. Un ciclo es una caminata de unos 2~$\mu$s medida con un reloj de 18.5 ns, de modo que lo que varía es, sobre todo, a qué lado de un flanco del reloj termina el ciclo. Por ejemplo, para $n_p = 3$, 35\,131 de 40\,976 ciclos tienen $\Delta t = 5$. La Tabla~\ref{tab:res_plugin} presenta la min-entropía condicional por estimación directa,
\begin{equation}
    H_\infty(\tau \mid n_p) = -\log_2 \sum_n \Pr[n_p = n]\, \max_t \Pr[\tau = t \mid n_p = n], \qquad \tau = \Delta t \bmod 256,
\end{equation}
bajo distintas condiciones de carga. La carga en los otros núcleos no cambia el resultado.

\begin{table}[H]
\centering
\caption{Min-entropía condicional por ciclo (estimación directa)}
\label{tab:res_plugin}
\begin{tabularx}{\textwidth}{@{}Xr@{}}
\toprule
\textbf{Condición} & \textbf{Bits por ciclo} \\
\midrule
Pi, en reposo, 5\,000\,000 de ciclos & 0.422 \\
Pi, fijada al núcleo 3, en reposo (2\,000\,000) & 0.437 \\
Pi, fijada al núcleo 3, núcleos 0--2 con carga (2\,000\,000) & 0.422 \\
Pi, sin fijar, núcleos 0--2 con carga (2\,000\,000) & 0.421 \\
Pi, condicionando además en 1 o 2 residuos previos & 0.406 / 0.404 \\
Computador con Windows (2\,000\,000) & 0.467 \\
\bottomrule
\end{tabularx}
\end{table}

La Tabla~\ref{tab:res_nist_ciclo} confirma estas cifras con \texttt{ea\_non\_iid}, aplicado al residuo $\Delta t - \operatorname{moda}(\Delta t \mid n_p) + 128$ de cada ciclo. Una entrada anterior, el residuo respecto de una única recta ajustada a todos los datos, obtuvo 0.565 en la Pi, pero es inválida: el desplazamiento fraccionario de la recta varía con $n_p$, de modo que el residuo redondeado todavía contiene información determinada por la semilla.

\begin{table}[H]
\centering
\caption{Min-entropía por ciclo según \texttt{ea\_non\_iid}}
\label{tab:res_nist_ciclo}
\small
\begin{tabular}{@{}lccccccc@{}}
\toprule
\textbf{Anfitrión} & \textbf{MCV} & \textbf{\textit{t-tuple}} & \textbf{MultiMMC} & \textbf{LZ78Y} & \textbf{$8 H_{\text{bits}}$} & \textbf{Estimado} & \textbf{Por byte} \\
\midrule
Raspberry Pi 4 & 0.419 & 0.408 & 0.419 & 0.420 & \textbf{0.375} & \textbf{0.375} & $\approx 1.9$ \\
Windows & 0.461 & 0.415 & 0.461 & 0.463 & \textbf{0.407} & \textbf{0.407} & $\approx 2.0$ \\
\bottomrule
\end{tabular}
\end{table}

Así, los anfitriones rápidos no son determinísticos, pero cada ciclo aporta entre 0.375 y 0.407 bits de min-entropía fresca según el NIST, es decir, unos 2 bits por byte de salida con la configuración por defecto, frente a los 8 bits por byte que contabiliza el diseño. Para respaldar 8 bits por byte a 0.375 bits por ciclo, la referencia necesitaría unos 22 ciclos por byte (unos 16 kB/s en la Pi), o bien la contabilidad explícita de la variante IoT aplicada al temporizador del anfitrión.

\paragraph{La batería completa no lo detecta.} La Tabla~\ref{tab:res_bateria_pi} presenta la batería completa sobre muestras de 1\,000\,000 de bytes de los cuatro candidatos de la librería, generadas en la Pi. SP 800-22 se ejecutó sobre los mismos archivos en WSL, con un directorio de trabajo por candidato y sin pruebas abortadas, y los valores de SP 800-90B coincidieron con los de la Pi en todos los dígitos.

\begin{table}[H]
\centering
\caption{Batería completa sobre 1\,000\,000 de bytes de los candidatos de la librería (Raspberry Pi 4)}
\label{tab:res_bateria_pi}
\begin{tabularx}{\textwidth}{@{}Yccccc@{}}
\toprule
\textbf{Candidato} & \textbf{Nivel 1} & \textbf{ENT $\chi^2$} & \textbf{\texttt{ea\_iid}} & \textbf{\texttt{ea\_non\_iid}} & \textbf{SP 800-22} \\
\midrule
Referencia \texttt{xorshift128+} & falla & 0.01\,\% & 7.85 & 7.42 & 1\,256/1\,296 \\
Referencia NEXT\_X48 & falla & 0.01\,\% & 7.85 & 7.19 & 1\,353/1\,374 \\
\texttt{xorshift128+} + SHA-256 & supera & 77.7\,\% & 7.89 & 7.23 & 1\,436/1\,452 \\
NEXT\_X48 + SHA-256 & supera & 78.8\,\% & 7.89 & 7.17 & 1\,332/1\,348 \\
\bottomrule
\end{tabularx}
\end{table}

Los totales de SP 800-22 difieren porque las pruebas \textit{RandomExcursions} solamente aplican a los flujos con suficientes ciclos. Los candidatos acondicionados fallan 16 combinaciones cada uno, frente a unas 14.5 y 13.5 esperadas por azar, y ninguna prueba falla en la proporción de flujos, salvo una plantilla de \textit{NonOverlappingTemplate} en el candidato NEXT\_X48 acondicionado (6 de 8 flujos), compatible con el azar entre 148 plantillas. Los crudos fallan justamente las pruebas que el sesgo del bit 7 debe romper: \textit{Frequency} y las dos \textit{CumulativeSums} en ambos, además de \textit{Runs} y \textit{Universal} en \texttt{xorshift128+}; la referencia \texttt{xorshift128+} falla 40 combinaciones, unas tres veces lo esperado por azar. Pero la conclusión central es otra: SP 800-90B asigna incluso a la salida \emph{cruda} entre 7.2 y 7.4 bits por byte, y la salida acondicionada supera todas las pruebas, aunque, según la evaluación por ciclo, contiene unos 1.9 bits de entropía fresca por byte. Estas herramientas miden cómo se ve la salida, no cuánta entrada impredecible la produjo. Solo la traza de la fuente revela el déficit, igual que solo la comparación entre reinicios reveló el determinismo de E1.

Durante este experimento se detectó y corrigió un error de procedimiento que vale la pena registrar. En la Pi, el ejecutable \texttt{assess} de STS se había instalado como un enlace simbólico; el arnés copia las plantillas de STS desde el directorio del ejecutable que encuentra en la ruta, que era el del enlace y no el de STS. Las plantillas no se copiaron, todas las ejecuciones de \textit{NonOverlappingTemplate} abortaron y el informe final contenía 148 filas residuales sin sentido, con totales de entre 1\,300 y 1\,450 combinaciones. El error se detectó porque la misma instalación no reproducía un informe conocido. Los resultados de la Tabla~\ref{tab:res_bateria_pi} provienen de la ejecución corregida, validada previamente contra ese informe.

\section{E5: validación de Ascon-AEAD128}
\label{sec:res_ascon}

La Tabla~\ref{tab:res_ascon_validacion} resume la validación de \texttt{ascon-iot}. Ninguna prueba utilizó QPP-RNG ni la aleatoriedad del sistema operativo: las entradas de las pruebas diferenciales se derivaron de forma determinística con SHAKE256.

\begin{table}[H]
\centering
\caption{Validación de \texttt{ascon-iot}}
\label{tab:res_ascon_validacion}
\begin{tabularx}{\textwidth}{@{}L{0.33\textwidth}Xr@{}}
\toprule
\textbf{Prueba} & \textbf{Detalle} & \textbf{Resultado} \\
\midrule
Vectores KAT & 1\,089 casos de Ascon-AEAD128, cifrado y descifrado; archivo KAT reproducido de forma independiente con \texttt{ascon-c} & 1\,089/1\,089 \\
ACVP (local, de extremo a extremo) & 60 casos de cifrado, 60 de descifrado (30 válidos aceptados, 30 inválidos rechazados), validados con GenVal & 120/120 \\
Diferencial contra \texttt{ascon-c} & Cifrado idéntico, y descifrado cruzado en ambas direcciones, con longitudes de borde hasta 1\,024 bytes & 3\,000/3\,000 \\
Autenticación negativa interna & Etiqueta, texto cifrado o datos asociados modificados; \textit{nonce} o clave incorrectos & 1\,280 rechazos, 0 aceptaciones \\
Autenticación negativa cruzada & Las mismas manipulaciones, enviadas a Rust y a C & 2\,560 rechazos, 0 aceptaciones \\
\bottomrule
\end{tabularx}
\end{table}

En total se ejecutaron 6\,840 verificaciones dirigidas, además de 768 descifrados válidos de control, sin discrepancias ni aceptaciones inesperadas. Estas pruebas aportan evidencia local y reproducible de equivalencia con la implementación de referencia y de rechazo ante fallas de autenticación, pero no constituyen una validación CAVP ni una certificación del NIST. La prueba de integración que usa QPP-RNG para generar la clave y el \textit{nonce} se superó.

\section{E6: desempeño y huella}
\label{sec:res_huella}

\subsection{Ascon-AEAD128 en el anfitrión}

La Tabla~\ref{tab:res_ascon_bench} presenta la línea base de desempeño de Ascon-AEAD128 en el computador portátil (Intel Core i5-1135G7, WSL2), sin datos asociados. Cada valor es la mediana de cinco ejecuciones independientes, y cada ejecución reporta la mediana de siete muestras, con el proceso fijado a un núcleo lógico.

\begin{table}[H]
\centering
\caption{Desempeño de Ascon-AEAD128 en el anfitrión (sin datos asociados)}
\label{tab:res_ascon_bench}
\begin{tabular}{@{}rrrrr@{}}
\toprule
\textbf{Mensaje} & \textbf{Cifrado (ns)} & \textbf{Descifrado (ns)} & \textbf{Cifrado (MiB/s)} & \textbf{Descifrado (MiB/s)} \\
\midrule
0 B & 121.57 & 129.01 & -- & -- \\
16 B & 165.72 & 170.25 & 92.1 & 89.6 \\
64 B & 294.04 & 289.03 & 207.6 & 211.2 \\
256 B & 822.16 & 740.37 & 296.9 & 329.8 \\
1\,024 B & 2\,875.23 & 2\,487.20 & 339.6 & 392.6 \\
4\,096 B & 11\,112.74 & 9\,469.40 & 351.5 & 412.5 \\
\bottomrule
\end{tabular}
\end{table}

El costo fijo por operación domina en los mensajes cortos, que son los habituales en IoT. Los datos asociados añaden un costo casi independiente del tamaño del mensaje, del orden de 42 a 44 ns por bloque de 16 bytes absorbido, y se observan saltos al cruzar los límites de 16 y 32 bytes, de acuerdo con la tasa de 16 bytes de Ascon-AEAD128. Construir el objeto de cifrado a partir de una clave ya disponible es despreciable frente a una operación AEAD. Estas cifras no son extrapolables a los microcontroladores.

\subsection{Huella de la pila completa en AVR}

La Tabla~\ref{tab:res_huella_avr} muestra el tamaño en Flash de la pila completa (fuente del perro guardián, \texttt{qpp-rng-iot}, SHA-256 y Ascon-AEAD128) en AVR, según los \textit{backends} de RustCrypto y el nivel de optimización. La imagen mínima genera la clave y el \textit{nonce} y cifra un mensaje; la segunda columna corresponde al \textit{firmware} de la prueba de extremo a extremo, que además descifra.

\begin{table}[H]
\centering
\caption{Flash de la pila completa en AVR}
\label{tab:res_huella_avr}
\small
\begin{tabularx}{\textwidth}{@{}>{\raggedright\arraybackslash}XL{0.27\textwidth}L{0.24\textwidth}@{}}
\toprule
\textbf{Compilación} & \textbf{Imagen mínima} & \textbf{Con descifrado} \\
\midrule
Por defecto (ATmega2560) & 136 KB (SHA-256: 79 KB; permutación de Ascon: 30 KB) & -- \\
+ \textit{backend} compacto de SHA-256 (ATmega2560) & 59 KB & -- \\
+ \textit{backend} compacto de Ascon, \texttt{opt-level = \textquotedbl{}s\textquotedbl{}} (ATmega328P) & 16.1 KB & 34.7 KB (no cabe) \\
+ \texttt{opt-level = \textquotedbl{}z\textquotedbl{}} (ATmega328P) & \textbf{15.7 KB} (48\,\%) & \textbf{21.0 KB} (64\,\%) \\
\bottomrule
\end{tabularx}
\end{table}

En RustCrypto, los \textit{backends} compactos se seleccionan con dos opciones de compilación: \texttt{sha2\_backend\_soft} para SHA-256 y \texttt{ascon\_backend} para Ascon. El camino de entropía de QPP-RNG ocupa menos de 3 KB. Con \texttt{opt-level = \textquotedbl{}s\textquotedbl{}}, el compilador replica la permutación de Ascon en cada función que la invoca (inicialización, datos asociados, finalización, cifrado y descifrado), con entre 3 y 5.5 KB por copia en AVR, de modo que una imagen que también descifra supera los 32 KB del ATmega328P. Con \texttt{opt-level = \textquotedbl{}z\textquotedbl{}}, la permutación queda en una sola función de 3 KB y ambas imágenes caben con holgura.

La RAM estática incluye las constantes del programa: en AVR, Rust las ubica en la sección \texttt{.data}, que se copia de la Flash a la RAM al arranque, entre ellas la tabla de 256 bytes de constantes de SHA-256. La imagen mínima usa así 366 bytes de RAM estática (298 de constantes y 68 de variables) y el \textit{firmware} de extremo a extremo, 662, por sus cadenas de texto del informe.

En el anfitrión, \texttt{cargo bloat} atribuye a \texttt{qpp-rng-reference} exactamente 732 bytes en los dos ejecutables con \texttt{xorshift128+} (crudo y acondicionado) y 628 bytes en los dos con NEXT\_X48, lo que confirma que el enlazador elimina el código de los candidatos que no se usan.

\subsection{Ejecución de extremo a extremo en las placas}
\label{sec:res_e2e}

Para verificar la integración en el dispositivo, se ejecutó la pila completa en las tres placas con un \textit{firmware} de prueba. Este construye \texttt{QppEntropy} sobre la fuente física de la placa, genera una clave de 128 bits y un \textit{nonce} inicial aleatorio, y luego cifra y descifra tres mensajes de 0, 16 y 64 bytes con 16 bytes de datos asociados, usando como \textit{nonce} de cada mensaje el inicial más su índice. Cada mensaje se cifra de nuevo con un bit alterado (en el texto cifrado o, en el mensaje vacío, en la etiqueta), y el descifrado debe rechazarlo. Un \textit{script} en el anfitrión reinicia la placa, lee el informe y recalcula cada texto cifrado y cada etiqueta con la implementación de referencia \texttt{ascon-c}, que antes se verificó contra los 1\,089 vectores KAT. El consumo de pila se midió pintando la RAM libre antes de ejecutar la tubería y buscando, al final, el byte más bajo sobrescrito, lo que incluye las interrupciones. Cada placa se reinició tres veces.

\begin{table}[H]
\centering
\caption{Ejecución de extremo a extremo de la pila completa (AVR a 16 MHz; nRF52840 a 64 MHz)}
\label{tab:res_e2e}
\small
\begin{tabularx}{\textwidth}{@{}>{\raggedright\arraybackslash}XL{0.2\textwidth}L{0.2\textwidth}L{0.2\textwidth}@{}}
\toprule
& \textbf{Nano (ATmega328P)} & \textbf{Mega 2560 (ATmega2560)} & \textbf{nRF52840} \\
\midrule
Reinicios con clave y \textit{nonce} distintos & 3 de 3 & 3 de 3 & 3 de 3 \\
Textos cifrados y etiquetas iguales a \texttt{ascon-c} & 9 de 9 & 9 de 9 & 9 de 9 \\
Descifrados correctos; alteraciones rechazadas & 9 de 9; 9 de 9 & 9 de 9; 9 de 9 & 9 de 9; 9 de 9 \\
Clave de 128 bits y su \textit{nonce} & 8.70 a 8.77~s & 8.82 a 8.85~s & 0.69 a 0.72~s \\
Cifrado de 0, 16 y 64 bytes (ms) & 4.38, 5.25 y 7.83 & 4.36, 5.23 y 7.81 & 0.13, 0.15 y 0.23 \\
Costo adicional por byte de mensaje & $\approx 860$ ciclos & $\approx 860$ ciclos & $\approx 103$ ciclos \\
Flash & 21\,016 B (64\,\%) & 21\,156 B (8\,\%) & 21\,664 B \\
RAM estática & 662 B & 662 B & 0 B \\
Pila máxima & 975 a 1\,029 B & 982 a 1\,039 B & 1\,032 B \\
RAM nunca usada & 357 B de 2\,048 (17\,\%) & 6\,491 B de 8\,192 & 261\,112 B de 262\,144 \\
\bottomrule
\end{tabularx}
\end{table}

Los tres reinicios de cada placa produjeron claves y \textit{nonces} distintos, los 27 textos cifrados y etiquetas coincidieron con \texttt{ascon-c}, todos los mensajes se descifraron correctamente y todas las alteraciones se rechazaron. El descifrado cuesta prácticamente lo mismo que el cifrado, a lo sumo un 3.1\,\% más. Con los \textit{backends} compactos, Ascon-AEAD128 tarda en los AVR 4.4~ms para un mensaje vacío con 16 bytes de datos asociados, más unos 860 ciclos por byte de mensaje, unas 34 veces más que en el nRF52840; los dos AVR coinciden dentro del 1\,\%, como corresponde al mismo núcleo a la misma frecuencia. En RAM, el Nano usa como máximo 1\,691 de sus 2\,048 bytes (662 estáticos y hasta 1\,029 de pila), de modo que a una aplicación real le quedan unos 350 bytes; en el nRF52840 el consumo es irrelevante frente a sus 256 KB.

\subsection{Tiempo para sembrar}

El acondicionador entrega su salida en bloques de 32 bytes, y cada bloque requiere 64 bytes crudos. Por eso, la primera salida tarda lo mismo sea una clave de 128 bits o una semilla de 256 bits: en la ejecución de extremo a extremo, la clave y su \textit{nonce}, que provienen del mismo bloque, estuvieron listos en 8.7~s en el Nano y 8.8~s en el Mega 2560, a unos 7.4 B/s de salida cruda, y en 0.7~s en el nRF52840. Cada bloque adicional de 32 bytes cuesta otros 64 bytes crudos, es decir, unos 3.7 B/s de salida acondicionada en AVR. Este tiempo de arranque es aceptable para generar la clave de un dispositivo, pero no para producir un \textit{nonce} aleatorio por mensaje; de ahí la recomendación de usar un contador o un DRBG sembrado al arranque.

\section{Síntesis de resultados}

La Tabla~\ref{tab:sintesis_resultados} resume los resultados frente a las preguntas de investigación.

\begin{table}[H]
\centering
\caption{Síntesis de resultados}
\label{tab:sintesis_resultados}
\begin{tabularx}{\textwidth}{@{}L{0.30\textwidth}Y@{}}
\toprule
\textbf{Pregunta} & \textbf{Resultado} \\
\midrule
¿Reproduce la referencia el comportamiento publicado? & Sí en lo estadístico: la salida cruda obtiene entre 7.2 y 7.4 bits por byte en SP 800-90B, y la acondicionada supera todas las pruebas. Además, se identificó un sesgo estructural del bit 7, que se predice con exactitud. \\
¿Hay entropía en el \textit{jitter} de ejecución en microcontroladores? & No. En las tres placas, la salida es idéntica bit a bit tras cada reinicio. En anfitriones rápidos sí la hay, pero solo unos 2 bits por byte, frente a los 8 contabilizados. \\
¿Qué fuente física la sustituye y a qué costo? & El perro guardián frente al Timer1 en AVR (1.62 y 2.24 bits por muestra; se acreditó 1) y el RNG en hardware del nRF52840 (7.24 bits por byte; se acreditaron 4). La variante IoT alcanza 7.4 B/s en AVR y 101 B/s en el nRF52840, con 2.9 a 6.5 KB de \textit{firmware}, y su salida supera la batería completa sobre 1\,000\,000 de bytes en las tres placas. \\
¿Se integra con Ascon-AEAD128 en un microcontrolador de 8 bits? & Sí. La pila completa ocupa entre 15.7 y 21.0 KB de Flash con los \textit{backends} compactos y cabe en un ATmega328P. En las tres placas generó una clave distinta en cada reinicio y cifró y descifró de forma idéntica a \texttt{ascon-c}; la primera clave tarda 8.7 a 8.8~s en AVR y 0.7~s en el nRF52840, y cifrar 64 bytes, 7.8 y 0.23~ms. Ascon supera KAT, ACVP y las pruebas diferenciales. \\
\bottomrule
\end{tabularx}
\end{table}
